\documentclass[11pt]{article}

\usepackage[a4paper,margin=1in]{geometry}
\usepackage{amsmath,amssymb,amsthm}
\usepackage{graphicx}
\usepackage[section]{placeins} % prevent floats from leaking into later sections
\graphicspath{{./}{./figs/}}
\usepackage{booktabs}
\usepackage{siunitx}
\usepackage{hyperref}
\usepackage[numbers]{natbib}

\title{Truels as Markov Games: Winner Landscapes and Exploitability of Heuristic and Learned Policies}
\author{Nazar~\textit{et al.}}
\date{\today}

\begin{document}
\maketitle

\begin{abstract}
A \emph{truel} is the three-player generalization of a duel: each player chooses a target (or abstains) until a single survivor remains.
Despite its simple rules, truels exhibit counterintuitive phenomena such as ``survival of the weakest,'' and their outcomes depend sharply on both marksmanship parameters and strategic choice.
We cast the truel as a finite general-sum Markov game and propose a reproducible evaluation protocol that turns ``interesting chaos'' into quantitative evidence:
(i) coarse parameter scans that produce winner heatmaps under multiple policy families, and
(ii) \emph{exploitability} heatmaps computed via best responses, measuring how far a policy profile is from equilibrium play.
We report full-grid winner regions for two classical baselines (``shoot the strongest'' and ``A abstains first'') and for a learned equilibrium-approximate profile (Nash-Q / Friend-or-Foe Q-learning).
Exploitability slices reveal where heuristic policies are robust and where they are systematically vulnerable.
Finally, we outline how truel-like three-party dynamics model realistic cybersecurity settings (attacker--defender--deceiver/insider), motivating equilibrium-oriented learning and vulnerability quantification.
\end{abstract}

\section{Introduction}

Truels have a weird charm: with three guns (or balloons, in the puzzle versions), ``being worse'' can be strategically better.
The phenomenon appears because the strongest player attracts attention, creating an implicit ``coalition of incentives'' against them even without explicit coordination.
As a result, the map from abilities to survival is not monotone.

\paragraph{History in brief.}
The truel entered the literature as a mathematical puzzle (often credited to Kinnaird's 1946 puzzle anthology) and later as an explicit game-theoretic object in the work of Shubik and others on ``does the fittest necessarily survive?'' \citep{Kinnaird1946,Shubik1954,Shubik1982}.
Formal game-theoretic analyses include Kilgour's early work on simultaneous and sequential variants \citep{Kilgour1972} and the widely-cited exposition by Brams and Kilgour \citep{BramsKilgour1997}.
Subsequent work explored solution concepts (cooperative vs.\ non-cooperative) \citep{BossertBramsKilgour2002}, Markov-chain analyses and winner distributions \citep{ToralAmengual2005,AmengualToral2006}, and behavioral/strategic variants such as abstention and ``shooting into the air'' \citep{Xu2012,DorrakiAllisonAbbott2019}.
Truels also appear in physics-flavored settings (e.g., quantum extensions) \citep{FlitneyAbbott2003}.
A consistent theme across these lines: the \emph{rules of play} and the \emph{strategies} matter as much as raw marksmanship.

\paragraph{Why bring reinforcement learning (RL) into this?}
Because the moment you say ``cybersecurity,'' you have an adaptive multi-agent system.
Attackers learn, defenders learn, and a third party is often present: an insider, a deception layer, or a competing attacker.
This naturally suggests \emph{stochastic games} \citep{Shapley1953}, aka Markov games in the RL literature \citep{Littman1994}.
But multi-agent RL is notoriously non-stationary: if you learn against a moving target, ``convergence'' becomes a mood rather than a theorem.
Equilibrium-based learning rules such as Nash-Q \citep{HuWellman2003} and structure-exploiting variants like Friend-or-Foe Q-learning \citep{Littman2001} offer a principled anchor: learn relative to equilibrium behavior rather than a particular opponent snapshot.

\paragraph{Our contribution.}
This paper is a methods-and-results bridge between classical truel analysis and modern evaluation in multi-agent learning.
We contribute:
\begin{itemize}
  \item A clean Markov-game formulation of a standard random truel with abstention.
  \item A coarse parameter scan over $(a,b,c)$ producing winner heatmaps under multiple policy families.
  \item An \emph{exploitability} evaluation (and heatmaps) over the same parameter space, using computed best responses as a worst-case vulnerability measure.
  \item A cybersecurity mapping that motivates why ``winner probabilities'' alone are not enough: we also care about how exploitable a deployed policy is.
\end{itemize}

\section{Background and Related Work}

\subsection{Truels and the ``survival of the weakest'' paradox}
The paradoxical regions where weaker players have higher win probability under rational play are reviewed and derived in several classical treatments \citep{BramsKilgour1997,AmengualToral2006}.
Markov-chain approaches make the dependence on rule variants explicit and provide closed-form win probabilities in common settings \citep{ToralAmengual2005,AmengualToral2006}.
Modern variants broaden the strategy set (abstention, suicidal/randomized behavior by weaker players) and study how these affect survival regions \citep{DorrakiAllisonAbbott2019}.
Philosophical and decision-theoretic angles also exist, e.g., analyzing when the worst shooter should abstain initially \citep{Xu2012}.

\subsection{Stochastic games, equilibrium, and multi-agent RL}
Stochastic games were introduced by Shapley \citep{Shapley1953} and provide the canonical model for sequential multi-agent interaction.
In RL, the Markov-game framing appears prominently in Littman's work \citep{Littman1994}.
Single-agent Q-learning is classical \citep{WatkinsDayan1992}; extending it to general-sum multi-agent settings typically requires a game-theoretic solution concept at each stage.
Nash-Q learning maintains joint-action value functions and computes Nash equilibria of stage games induced by the Q-values \citep{HuWellman2003}.
Friend-or-Foe Q-learning simplifies the equilibrium computation when opponents are known to be either aligned (friends) or strictly adversarial (foes), retaining convergence guarantees in those cases \citep{Littman2001}.
For bounded rationality or smooth best responses, quantal response equilibrium provides another link between optimization and realistic behavior \citep{McKelveyPalfrey1995}.

\subsection{Exploitability and vulnerability metrics}
For deployment, a policy is only as good as its worst-case exposure.
\emph{Exploitability} measures the gain an adversary can achieve by best-responding to a policy profile; in two-player zero-sum games it aligns cleanly with distance to Nash equilibrium.
In large games, exact computation can be hard, motivating approximate best responses and approximate exploitability \citep{TimbersEtAl2020}.
Regret-minimization methods such as Counterfactual Regret Minimization (CFR) are standard for approaching Nash equilibria and connect directly to exploitability in extensive-form games \citep{ZinkevichEtAl2007}.

\section{Model: A Truel as a Finite Markov Game}

We consider three players $A,B,C$ with marksmanship parameters $a,b,c \in (0,1)$, interpreted as hit probabilities when they choose a target.
At any time, the \emph{state} is the set of surviving players:
\[
\mathcal{S}=\{\{A,B,C\},\{A,B\},\{A,C\},\{B,C\},\{A\},\{B\},\{C\}\}.
\]
Absorbing terminal states are $\{A\},\{B\},\{C\}$.

\paragraph{Turn selection (random truel).}
At each step, one of the currently alive players is chosen uniformly at random to act \citep{AmengualToral2006}.

\paragraph{Actions.}
If player $i$ is chosen and at least one opponent is alive, $i$ selects an action
\[
u_i \in \{\text{shoot }j \ \text{for each alive } j\neq i\}\ \cup\ \{\text{abstain}\}.
\]
If $i$ shoots $j$, then $j$ is removed from the state with probability equal to $i$'s marksmanship; otherwise the state remains unchanged.
If $i$ abstains, the state remains unchanged.

\paragraph{Utilities.}
Each player receives terminal utility $1$ if they are the sole survivor and $0$ otherwise.
Equivalently, expected utility equals the survival (win) probability from the initial state.

This defines a finite episodic general-sum Markov game.

\section{Policies, Learning, and Evaluation}

\subsection{Policy families}
We study three profiles:
\begin{enumerate}
  \item \textbf{Baseline Strongest (BS):} each acting player shoots the strongest remaining opponent (ties broken deterministically).
  \item \textbf{Baseline A-pass (AP):} player $A$ abstains in the initial $ABC$ state; otherwise players follow the strongest-opponent rule.
  \item \textbf{Learned equilibrium-approx (LQ):} a learned profile produced by Nash-Q \citep{HuWellman2003} or Friend-or-Foe Q-learning \citep{Littman2001}, trained on the Markov-game dynamics above.
\end{enumerate}

\subsection{Coarse parameter scan and winner heatmaps}
We scan a coarse grid $a,b,c \in \{0.1,0.2,\dots,0.9\}$ and, for each point, estimate (or compute exactly, since the state space is tiny) the win probabilities under each policy profile.
We then color each grid point by $\arg\max\{\pi_A,\pi_B,\pi_C\}$ to obtain winner heatmaps.

\subsection{Exploitability}
Fix a policy profile $\pi=(\pi_A,\pi_B,\pi_C)$ and let $u_i(\pi)$ be player $i$'s expected payoff (win probability) from the start state.
Define player $i$'s \emph{deviation gain} (one-player best response) by
\[
\delta_i(\pi)= \max_{\pi_i'} u_i(\pi_i',\pi_{-i}) - u_i(\pi).
\]
The \emph{(one-step) exploitability} of the profile is then
\[
\mathrm{Exploit}(\pi)=\sum_{i\in\{A,B,C\}}\delta_i(\pi),
\]
and we also report the vector $(\delta_A,\delta_B,\delta_C)$.
This is the ``how much can a rational attacker improve by deviating'' metric: if it is near zero, the profile is close to a Nash equilibrium in the sense of $\epsilon$-Nash \citep{TimbersEtAl2020}.

\paragraph{Computation in this truel.}
Because $\mathcal{S}$ is tiny and actions are discrete, best responses can be computed by dynamic programming / value iteration on the induced MDP for each deviating player (holding the other two fixed).
This makes \emph{full-grid exploitability maps} feasible.

\section{Results}

\subsection{Paired winner--exploitability heatmaps}
We visualize (i) the winning label (argmax win probability) and (ii) total exploitability across $(b,c)$ for three $a$-slices. These paired plots separate performance from strategic robustness.

\begin{figure}[!htbp]
  \centering
  \begin{minipage}[t]{0.48\linewidth}
    \centering\includegraphics[width=\linewidth]{\detokenize{figs/winner_policy_shoot_strongest_a_0.1.pdf}}
  \end{minipage}
  \begin{minipage}[t]{0.48\linewidth}
    \centering\includegraphics[width=\linewidth]{\detokenize{figs/exploit_policy_shoot_strongest_a_0.1.pdf}}
  \end{minipage}
  \caption{Paired heatmaps for policy \texttt{shoot\_strongest} at $a=0.1$. Left: winner label (0=A,1=B,2=C) computed from exact absorption probabilities. Right: total exploitability $\sum_i \mathrm{Exploit}_i$ (payoff scale).}
  \label{fig:shoot-strongest:a0p1}
\end{figure}

\begin{figure}[!htbp]
  \centering
  \begin{minipage}[t]{0.48\linewidth}
    \centering\includegraphics[width=\linewidth]{\detokenize{figs/winner_policy_shoot_strongest_a_0.5.pdf}}
  \end{minipage}
  \begin{minipage}[t]{0.48\linewidth}
    \centering\includegraphics[width=\linewidth]{\detokenize{figs/exploit_policy_shoot_strongest_a_0.5.pdf}}
  \end{minipage}
  \caption{Paired heatmaps for policy \texttt{shoot\_strongest} at $a=0.5$. Left: winner label (0=A,1=B,2=C) computed from exact absorption probabilities. Right: total exploitability $\sum_i \mathrm{Exploit}_i$ (payoff scale).}
  \label{fig:shoot-strongest:a0p5}
\end{figure}

\begin{figure}[!htbp]
  \centering
  \begin{minipage}[t]{0.48\linewidth}
    \centering\includegraphics[width=\linewidth]{\detokenize{figs/winner_policy_shoot_strongest_a_0.9.pdf}}
  \end{minipage}
  \begin{minipage}[t]{0.48\linewidth}
    \centering\includegraphics[width=\linewidth]{\detokenize{figs/exploit_policy_shoot_strongest_a_0.9.pdf}}
  \end{minipage}
  \caption{Paired heatmaps for policy \texttt{shoot\_strongest} at $a=0.9$. Left: winner label (0=A,1=B,2=C) computed from exact absorption probabilities. Right: total exploitability $\sum_i \mathrm{Exploit}_i$ (payoff scale).}
  \label{fig:shoot-strongest:a0p9}
\end{figure}

\begin{figure}[!htbp]
  \centering
  \begin{minipage}[t]{0.48\linewidth}
    \centering\includegraphics[width=\linewidth]{\detokenize{figs/winner_policy_A_pass_a_0.1.pdf}}
  \end{minipage}
  \begin{minipage}[t]{0.48\linewidth}
    \centering\includegraphics[width=\linewidth]{\detokenize{figs/exploit_policy_A_pass_a_0.1.pdf}}
  \end{minipage}
  \caption{Paired heatmaps for policy \texttt{A\_pass} at $a=0.1$. Left: winner label (0=A,1=B,2=C) computed from exact absorption probabilities. Right: total exploitability $\sum_i \mathrm{Exploit}_i$ (payoff scale).}
  \label{fig:A-pass:a0p1}
\end{figure}

\begin{figure}[!htbp]
  \centering
  \begin{minipage}[t]{0.48\linewidth}
    \centering\includegraphics[width=\linewidth]{\detokenize{figs/winner_policy_A_pass_a_0.5.pdf}}
  \end{minipage}
  \begin{minipage}[t]{0.48\linewidth}
    \centering\includegraphics[width=\linewidth]{\detokenize{figs/exploit_policy_A_pass_a_0.5.pdf}}
  \end{minipage}
  \caption{Paired heatmaps for policy \texttt{A\_pass} at $a=0.5$. Left: winner label (0=A,1=B,2=C) computed from exact absorption probabilities. Right: total exploitability $\sum_i \mathrm{Exploit}_i$ (payoff scale).}
  \label{fig:A-pass:a0p5}
\end{figure}

\begin{figure}[!htbp]
  \centering
  \begin{minipage}[t]{0.48\linewidth}
    \centering\includegraphics[width=\linewidth]{\detokenize{figs/winner_policy_A_pass_a_0.9.pdf}}
  \end{minipage}
  \begin{minipage}[t]{0.48\linewidth}
    \centering\includegraphics[width=\linewidth]{\detokenize{figs/exploit_policy_A_pass_a_0.9.pdf}}
  \end{minipage}
  \caption{Paired heatmaps for policy \texttt{A\_pass} at $a=0.9$. Left: winner label (0=A,1=B,2=C) computed from exact absorption probabilities. Right: total exploitability $\sum_i \mathrm{Exploit}_i$ (payoff scale).}
  \label{fig:A-pass:a0p9}
\end{figure}

\begin{figure}[!htbp]
  \centering
  \begin{minipage}[t]{0.48\linewidth}
    \centering\includegraphics[width=\linewidth]{\detokenize{figs/winner_policy_nash_approx_a_0.1.pdf}}
  \end{minipage}
  \begin{minipage}[t]{0.48\linewidth}
    \centering\includegraphics[width=\linewidth]{\detokenize{figs/exploit_policy_nash_approx_a_0.1.pdf}}
  \end{minipage}
  \caption{Paired heatmaps for policy \texttt{nash\_approx} at $a=0.1$. Left: winner label (0=A,1=B,2=C) computed from exact absorption probabilities. Right: total exploitability $\sum_i \mathrm{Exploit}_i$ (payoff scale).}
  \label{fig:nash-approx:a0p1}
\end{figure}

\begin{figure}[!htbp]
  \centering
  \begin{minipage}[t]{0.48\linewidth}
    \centering\includegraphics[width=\linewidth]{\detokenize{figs/winner_policy_nash_approx_a_0.5.pdf}}
  \end{minipage}
  \begin{minipage}[t]{0.48\linewidth}
    \centering\includegraphics[width=\linewidth]{\detokenize{figs/exploit_policy_nash_approx_a_0.5.pdf}}
  \end{minipage}
  \caption{Paired heatmaps for policy \texttt{nash\_approx} at $a=0.5$. Left: winner label (0=A,1=B,2=C) computed from exact absorption probabilities. Right: total exploitability $\sum_i \mathrm{Exploit}_i$ (payoff scale).}
  \label{fig:nash-approx:a0p5}
\end{figure}

\begin{figure}[!htbp]
  \centering
  \begin{minipage}[t]{0.48\linewidth}
    \centering\includegraphics[width=\linewidth]{\detokenize{figs/winner_policy_nash_approx_a_0.9.pdf}}
  \end{minipage}
  \begin{minipage}[t]{0.48\linewidth}
    \centering\includegraphics[width=\linewidth]{\detokenize{figs/exploit_policy_nash_approx_a_0.9.pdf}}
  \end{minipage}
  \caption{Paired heatmaps for policy \texttt{nash\_approx} at $a=0.9$. Left: winner label (0=A,1=B,2=C) computed from exact absorption probabilities. Right: total exploitability $\sum_i \mathrm{Exploit}_i$ (payoff scale).}
  \label{fig:nash-approx:a0p9}
\end{figure}




We include:
(i) winner heatmaps for each policy family, and
(ii) exploitability heatmaps (and slices at fixed $a$ values) for the learned and baseline profiles.

\paragraph{Winner landscapes.}
Baseline Strongest reproduces the known non-monotonic winner regions and ``survival of the weakest'' pockets observed in classical analyses \citep{AmengualToral2006,DorrakiAllisonAbbott2019}.
Baseline A-pass shifts these regions and, in some $(a,b,c)$ settings, improves $A$'s survival probability by avoiding being the first ``magnet'' for incoming fire.

\paragraph{Exploitability landscapes.}
Exploitability heatmaps identify parameter regions where a heuristic profile is brittle:
even if it yields a decent win probability in expectation, a single best-responding opponent can substantially increase their chance of winning.
The learned equilibrium-approx profile typically reduces exploitability, but the reduction is not uniform across the parameter space; the slices show where learning remains vulnerable and where it becomes robust.

\paragraph{A useful plot: winner vs.\ exploitability.}
Combining winner identity with exploitability (e.g., scatter or binned maps) reveals cases where a policy ``wins'' but remains highly exploitable—exactly the kind of trap that matters in adversarial domains.

\section{Cybersecurity Interpretation}

Truels are not literally gunfights in practice (sadly, most networks are made of packets, not pistols), but the structure is familiar:

\begin{itemize}
  \item \textbf{Attacker--Defender--Insider:} an external attacker ($C$), a defender ($A$), and an insider threat or compromised component ($B$).
  \item \textbf{Attacker--Defender--Deception:} defense deploys deception (honeypots, decoys) that behaves like a third strategic entity influencing targeting incentives.
  \item \textbf{Competing attackers:} two attackers compete for access and resources while a defender adapts.
\end{itemize}

In such settings, ``best shot'' is a stand-in for capability (exploit reliability, detection quality, response speed), and the key operational question is not only ``who wins'' but ``how exploitable is our deployed policy?''
Exploitability maps provide a vulnerability landscape: where small deviations (new attacker tactics) can extract large gains, a system is strategically fragile.

\section{Conclusion and Next Steps}

We presented a compact but evidence-driven pipeline for studying truels as Markov games:
winner heatmaps plus exploitability heatmaps.
This connects classical paradoxes in truels to a modern robustness metric used in game AI and adversarial learning.

Near-term extensions:
(i) finer grids and confidence bounds,
(ii) richer action sets (costly shots, imperfect observation, limited ammo),
(iii) moving from one-player deviations to coalition deviations,
(iv) explicit cybersecurity toy environments where $a,b,c$ emerge from measurable system properties.

\section*{Reproducibility}
We provide scripts to reproduce all parameter scans and exploitability maps, and we release the aggregated CSV outputs used to generate the figures.


\FloatBarrier
\clearpage
\appendix
\section{Per-player exploitability maps}
For $a=0.1$ we additionally report per-player exploitability maps. Each panel shows the payoff improvement obtainable by a single player deviating to a best response.

\begin{figure}[!htbp]
  \centering
  \begin{minipage}[t]{0.32\linewidth}\centering\includegraphics[width=\linewidth]{\detokenize{figs/exploit_A_policy_shoot_strongest_a_0.1.pdf}}\end{minipage}
  \begin{minipage}[t]{0.32\linewidth}\centering\includegraphics[width=\linewidth]{\detokenize{figs/exploit_B_policy_shoot_strongest_a_0.1.pdf}}\end{minipage}
  \begin{minipage}[t]{0.32\linewidth}\centering\includegraphics[width=\linewidth]{\detokenize{figs/exploit_C_policy_shoot_strongest_a_0.1.pdf}}\end{minipage}
  \caption{Per-player exploitability at $a=0.1$ for policy \texttt{shoot\_strongest}. Left: player A; middle: player B; right: player C.}
  \label{fig:perplayer:shoot-strongest:a0p1}
\end{figure}

\begin{figure}[!htbp]
  \centering
  \begin{minipage}[t]{0.32\linewidth}\centering\includegraphics[width=\linewidth]{\detokenize{figs/exploit_A_policy_A_pass_a_0.1.pdf}}\end{minipage}
  \begin{minipage}[t]{0.32\linewidth}\centering\includegraphics[width=\linewidth]{\detokenize{figs/exploit_B_policy_A_pass_a_0.1.pdf}}\end{minipage}
  \begin{minipage}[t]{0.32\linewidth}\centering\includegraphics[width=\linewidth]{\detokenize{figs/exploit_C_policy_A_pass_a_0.1.pdf}}\end{minipage}
  \caption{Per-player exploitability at $a=0.1$ for policy \texttt{A\_pass}. Left: player A; middle: player B; right: player C.}
  \label{fig:perplayer:A-pass:a0p1}
\end{figure}

\begin{figure}[!htbp]
  \centering
  \begin{minipage}[t]{0.32\linewidth}\centering\includegraphics[width=\linewidth]{\detokenize{figs/exploit_A_policy_nash_approx_a_0.1.pdf}}\end{minipage}
  \begin{minipage}[t]{0.32\linewidth}\centering\includegraphics[width=\linewidth]{\detokenize{figs/exploit_B_policy_nash_approx_a_0.1.pdf}}\end{minipage}
  \begin{minipage}[t]{0.32\linewidth}\centering\includegraphics[width=\linewidth]{\detokenize{figs/exploit_C_policy_nash_approx_a_0.1.pdf}}\end{minipage}
  \caption{Per-player exploitability at $a=0.1$ for policy \texttt{nash\_approx}. Left: player A; middle: player B; right: player C.}
  \label{fig:perplayer:nash-approx:a0p1}
\end{figure}


\FloatBarrier
\clearpage
\bibliographystyle{plainnat}
\bibliography{truel_refs}

\end{document}
